\section{Boolean M\"{o}bius inversion}

\subsection{Alternating sums over supersets}

\begin{lemma}[Alternating sum over supersets]
\label{lem.pie.two-sets-altsum}
\lean{AlgebraicCombinatorics.InclusionExclusion.alternating_sum_supersets}
\leanok
Let $P \subseteq Q$ be two finite sets.

\textbf{(a)} We have
\begin{equation}
\sum_{\substack{I\subseteq Q;\\P\subseteq I}}\left(  -1\right)  ^{\left\vert
I\right\vert }=\left(  -1\right)  ^{\left\vert P\right\vert }\left[
P=Q\right]  . \label{eq.lem.pie.two-sets-altsum.a}
\end{equation}


\textbf{(b)} We have
\begin{equation}
\sum_{\substack{I\subseteq Q;\\P\subseteq I}}\left(  -1\right)  ^{\left\vert
Q\setminus I\right\vert }=\left[  P=Q\right]  .
\label{eq.lem.pie.two-sets-altsum.b}
\end{equation}
\end{lemma}

\begin{proof}
\textbf{(a)} We must prove the equality (\ref{eq.lem.pie.two-sets-altsum.a}). If $P=Q$,
then the only subset $I$ of $Q$ satisfying $P\subseteq I$ is $Q$ itself,
so the sum equals $(-1)^{|Q|} = (-1)^{|P|} \cdot [P = Q]$
(since $[Q = Q] = 1$).

For the rest of this proof, we assume $P\neq Q$. Thus, $\left[  P=Q\right]  =0$.

Now, $P$ is a proper subset of $Q$, so there exists some $q\in Q$
such that $q\notin P$. Fix such a $q$.

Let
\[
\mathcal{A}:=\left\{  I\subseteq Q\ \mid\ P\subseteq I\right\},
\]
and let $\operatorname{sign} I := (-1)^{|I|}$ for each $I\in\mathcal{A}$.
Then
\begin{equation}
\sum_{I\in\mathcal{A}}\operatorname{sign} I=\sum_{\substack{I\subseteq
Q;\\P\subseteq I}}\left(  -1\right)  ^{\left\vert I\right\vert }.
\label{pf.lem.pie.two-sets-altsum.a.3}
\end{equation}
We define a map $f:\mathcal{A}\rightarrow\mathcal{A}$ by
\[
f\left(  I\right)  :=I\bigtriangleup\left\{  q\right\}  =
\begin{cases}
I\setminus\left\{  q\right\}  , & \text{if }q\in I;\\
I\cup\left\{  q\right\}  , & \text{if }q\notin I
\end{cases}
\ \ \ \ \ \ \ \ \ \ \text{for each }I\in\mathcal{A}.
\]
This map $f$ is well-defined
and it is an involution (the symmetric difference with $\{q\}$ undoes itself).
This involution has no fixed points
(since $f(I) = I \bigtriangleup \{q\} \neq I$).
Furthermore, for each $I \in \mathcal{A}$, the set $f(I) = I \bigtriangleup \{q\}$
differs from $I$ in exactly one element (namely $q$), so
$|f(I)| = |I| \pm 1$. Therefore
\[
(-1)^{|f(I)|} = -(-1)^{|I|},
\]
i.e., $\operatorname{sign}(f(I)) = -\operatorname{sign} I$.
Since $f$ is a sign-reversing fixed-point-free involution on $\mathcal{A}$,
the elements of $\mathcal{A}$ pair up into pairs with opposite signs, giving
\[
\sum_{I\in\mathcal{A}}\operatorname{sign} I = 0.
\]
Comparing with (\ref{pf.lem.pie.two-sets-altsum.a.3}), we obtain
\[
\sum_{\substack{I\subseteq Q;\\P\subseteq I}}\left(  -1\right)  ^{\left\vert
I\right\vert }=0=\left(  -1\right)  ^{\left\vert P\right\vert }\left[
P=Q\right]
\]
(since $(-1)^{|P|} \cdot 0 = 0$).
This proves (\ref{eq.lem.pie.two-sets-altsum.a}).
\medskip

\textbf{(b)} If $I$ is any subset of $Q$, then $\left\vert Q\setminus
I\right\vert =\left\vert Q\right\vert -\left\vert I\right\vert \equiv
\left\vert Q\right\vert +\left\vert I\right\vert \pmod{2}$, so
\[
\left(  -1\right)  ^{\left\vert Q\setminus I\right\vert }=\left(  -1\right)
^{\left\vert Q\right\vert +\left\vert I\right\vert }=\left(  -1\right)
^{\left\vert Q\right\vert }\left(  -1\right)  ^{\left\vert I\right\vert }
\]
(since $\left\vert Q\setminus I\right\vert \equiv
\left\vert Q\right\vert +\left\vert I\right\vert \pmod{2}$).
Hence
\begin{align*}
\sum_{\substack{I\subseteq Q;\\P\subseteq I}}\left(  -1\right)
^{\left\vert Q\setminus I\right\vert }
&=\sum_{\substack{I\subseteq Q;\\P\subseteq I}}\left(  -1\right)  ^{\left\vert
Q\right\vert }\left(  -1\right)  ^{\left\vert I\right\vert }
=\left(-1\right)  ^{\left\vert Q\right\vert }\underbrace{\sum_{\substack{I\subseteq
Q;\\P\subseteq I}}\left(  -1\right)  ^{\left\vert I\right\vert }%
}_{\substack{=\left(  -1\right)  ^{\left\vert P\right\vert }\left[
P=Q\right]  \\\text{(by part \textbf{(a)})}}}\\
&=\left(  -1\right)  ^{\left\vert Q\right\vert }\left(  -1\right)
^{\left\vert P\right\vert }\left[  P=Q\right]  .
\end{align*}
It remains to check that
\begin{equation}
\left(  -1\right)  ^{\left\vert Q\right\vert }\left(  -1\right)  ^{\left\vert
P\right\vert }\left[  P=Q\right]  =\left[  P=Q\right].
\label{pf.lem.pie.two-sets-altsum.b.1}
\end{equation}
If $P\neq Q$, then $[P = Q] = 0$ and both sides are $0$.
If $P = Q$, then $(-1)^{|Q|}(-1)^{|P|} = (-1)^{2|Q|} = 1$,
so both sides equal $1$.
Thus
\[
\sum_{\substack{I\subseteq Q;\\P\subseteq I}}\left(  -1\right)  ^{\left\vert
Q\setminus I\right\vert }=\left[  P=Q\right]  .
\]
This proves (\ref{eq.lem.pie.two-sets-altsum.b}).
\end{proof}

\subsection{The Boolean M\"{o}bius inversion formula}

\begin{lemma}
\label{lem.pie.sum-swap}
\lean{BooleanMobius.sum_powerset_powerset_swap}
\leanhelper
\leanok
Let $Q$ be a finite set, $A$ an additive abelian group, and
$f$ a function that assigns to each pair $(J, P)$ of subsets of $Q$ an
element of $A$. Then
\[
\sum_{J \subseteq Q} \sum_{P \subseteq J} f(J, P)
= \sum_{P \subseteq Q} \sum_{\substack{J \subseteq Q;\\P \subseteq J}} f(J, P).
\]
\end{lemma}

\begin{proof}
Both sides sum the same function $f$ over all pairs $(J, P)$ of subsets
of $Q$ satisfying $P \subseteq J \subseteq Q$; only the order of
summation differs.
\end{proof}

\begin{theorem}
[Boolean M\"{o}bius inversion]\label{thm.pie.moeb}
\lean{BooleanMobius.booleanMobiusInversion}
\leantarget
\leanok
Let $S$ be a finite set. Let
$A$ be any additive abelian group.

For each subset $I$ of $S$, let $a_{I}$ and $b_{I}$ be two elements of $A$.

Assume that
\begin{equation}
b_{I}=\sum_{J\subseteq I}a_{J}\ \ \ \ \ \ \ \ \ \ \text{for all }I\subseteq S.
\label{eq.thm.pie.moeb.ass}
\end{equation}


Then, we also have
\begin{equation}
a_{I}=\sum_{J\subseteq I}\left(  -1\right)  ^{\left\vert I\setminus
J\right\vert }b_{J}\ \ \ \ \ \ \ \ \ \ \text{for all }I\subseteq S.
\label{eq.thm.pie.moeb.claim}
\end{equation}
\end{theorem}

\begin{proof}
Fix a subset $Q$ of $S$. We shall prove that
\[
a_{Q}=\sum_{I\subseteq Q}\left(  -1\right)  ^{\left\vert Q\setminus
I\right\vert }b_{I}.
\]


We begin by rewriting the right hand side.
By (\ref{eq.thm.pie.moeb.ass}), we have $b_I = \sum_{J \subseteq I} a_J$ for each $I$, so
\begin{align}
\sum_{I\subseteq Q}\left(  -1\right)  ^{\left\vert Q\setminus I\right\vert
}b_{I}  &  =\sum_{I\subseteq Q}\left(  -1\right)
^{\left\vert Q\setminus I\right\vert }\underbrace{\sum_{J\subseteq I}a_{J}%
}_{=\sum_{P\subseteq I}a_{P}}=\sum_{I\subseteq Q}\left(  -1\right)
^{\left\vert Q\setminus I\right\vert }\sum_{P\subseteq I}a_{P}\nonumber\\
&  =\sum_{I\subseteq Q}\ \ \sum_{P\subseteq I}\left(  -1\right)  ^{\left\vert
Q\setminus I\right\vert }a_{P}. \label{pf.thm.pie.moeb.1}
\end{align}


The two summation signs ``$\sum_{I\subseteq Q}\ \ \sum
_{P\subseteq I}$'' on the right hand side of this equality
result in a sum over all pairs $\left(  I,P\right)  $ of subsets of $Q$
satisfying $P\subseteq I\subseteq Q$. The same result can be obtained by the
two summation signs ``$\sum_{P\subseteq Q}\ \ \sum
_{\substack{I\subseteq Q;\\P\subseteq I}}$''
(by Lemma~\ref{lem.pie.sum-swap}).
Hence, (\ref{pf.thm.pie.moeb.1})
rewrites as follows:
\begin{align}
\sum_{I\subseteq Q}\left(  -1\right)  ^{\left\vert Q\setminus I\right\vert
}b_{I}  &  =\sum_{P\subseteq Q}\ \ \sum_{\substack{I\subseteq Q;\\P\subseteq
I}}\left(  -1\right)  ^{\left\vert Q\setminus I\right\vert }a_{P}\nonumber\\
&  =\sum_{P\subseteq Q}\left(  \sum_{\substack{I\subseteq Q;\\P\subseteq
I}}\left(  -1\right)  ^{\left\vert Q\setminus I\right\vert }\right)  a_{P}.
\label{pf.thm.pie.moeb.2}
\end{align}


We want to prove that this equals $a_{Q}$. To this end, we
show that the coefficient $\sum_{\substack{I\subseteq
Q;\\P\subseteq I}}\left(  -1\right)  ^{\left\vert Q\setminus I\right\vert }$
in front of $a_{P}$ is $0$ whenever $P\neq Q$, and is $1$ whenever $P=Q$.
That is, every subset $P$ of $Q$
satisfies
\begin{equation}
\sum_{\substack{I\subseteq Q;\\P\subseteq I}}\left(  -1\right)  ^{\left\vert
Q\setminus I\right\vert }=\left[  P=Q\right]  . \label{pf.thm.pie.moeb.4}
\end{equation}

This is precisely Lemma \ref{lem.pie.two-sets-altsum} \textbf{(b)}.

Now, (\ref{pf.thm.pie.moeb.2}) becomes
(using (\ref{pf.thm.pie.moeb.4}))
\begin{align*}
\sum_{I\subseteq Q}\left(  -1\right)  ^{\left\vert Q\setminus I\right\vert
}b_{I}  &  =\sum_{P\subseteq Q}\underbrace{\left(  \sum_{\substack{I\subseteq
Q;\\P\subseteq I}}\left(  -1\right)  ^{\left\vert Q\setminus I\right\vert
}\right)  }_{=\left[  P=Q\right]  }a_{P}=\sum_{P\subseteq Q}\left[  P=Q\right]
a_{P}\\
&  =\underbrace{\left[  Q=Q\right]  }_{\substack{=1\\\text{(since }%
Q=Q\text{)}}}a_{Q}+\sum_{\substack{P\subseteq Q;\\P\neq Q}}\underbrace{\left[
P=Q\right]  }_{\substack{=0\\\text{(since }P\neq Q\text{)}}}a_{P}\\
&  =a_{Q}+\underbrace{\sum_{\substack{P\subseteq Q;\\P\neq Q}}0 \cdot a_{P}}%
_{=0}=a_{Q}.
\end{align*}
In other words, $a_{Q}=\sum_{I\subseteq Q}\left(  -1\right)  ^{\left\vert
Q\setminus I\right\vert }b_{I}$.

Since $Q$ was an arbitrary subset of $S$, we have shown that
\[
a_{I}=\sum_{J\subseteq I}\left(  -1\right)  ^{\left\vert I\setminus
J\right\vert }b_{J}\ \ \ \ \ \ \ \ \ \ \text{for all }I\subseteq S.
\]
This proves Theorem \ref{thm.pie.moeb}.
\end{proof}

\begin{lemma}
\label{lem.pie.moeb-module}
\lean{BooleanMobius.booleanMobiusInversion'}
\leanhelper
\leanok
Alternative formulation of Boolean M\"{o}bius inversion
(Theorem~\ref{thm.pie.moeb}) using the $\mathbb{Z}$-module scalar
multiplication: under the same hypotheses, for all $I\subseteq S$,
\[
a_{I}=\sum_{J\subseteq I}\left(-1\right)^{\left\vert I\setminus
J\right\vert}\cdot b_{J},
\]
where the multiplication $(-1)^{|I\setminus J|}\cdot b_J$ denotes the
$\mathbb{Z}$-module scalar action.
\end{lemma}

\begin{proof}
\leanok
Immediate from Theorem~\ref{thm.pie.moeb}, since the integer scalar
multiplication $n\cdot a$ in any additive abelian group coincides with
$n$ times $a$ when the group carries a $\mathbb{Z}$-module structure.
\end{proof}
